\documentclass[conference]{IEEEtran}
\IEEEoverridecommandlockouts

\usepackage{cite}
\usepackage{amsmath,amssymb}
\usepackage{graphicx}
\usepackage{booktabs}
\usepackage{url}
\usepackage[hidelinks]{hyperref}

\begin{document}

\title{Batch-Tiled Evaluation Overstates the Cost of\\Causal Sleep Staging}

\author{\IEEEauthorblockN{Jahnvi R, Swapnil S, Shreeya S Methuku, Shreshtha Dixit}
\IEEEauthorblockA{Department of Computer Science (AIML) and Engineering,
PES University, Bengaluru, India}}

\maketitle

\begin{abstract}
A sleep-staging model intended for a wearable device may use only past signal, and this
constraint is generally assumed to cost accuracy. We measure that cost with a matched
ablation in which one architecture is trained with and without access to future signal,
holding parameter count, data, folds and seeds fixed, and find that the answer depends
entirely on how the model is evaluated. Under the batch-tiled protocol standard in this
literature, which covers a recording with non-overlapping context windows, the causal
constraint costs $0.025$ Cohen's $\kappa$ on 78 Sleep-EDF subjects and appears concentrated
in REM. We show this is an artifact. Tiling scores predictions at every position of the
context buffer, including positions with almost no accumulated history, which a streaming
system passes through once at start-up; the causal model is crippled there while the
non-causal model is not, because it retains lookahead. The tiled score is arithmetically the
mean of a curve over buffer positions, which we verify to within $0.003\,\kappa$ for both
arms, and the entire penalty is generated by positions with under roughly 60 seconds of
history. Evaluated at matched latency the penalty disappears: the causal model reaches
$\kappa = 0.700$ with zero delay while its non-causal twin does not exceed $\kappa = 0.698$
at any delay within the buffer both were trained for, and the apparent REM concentration
inverts. The practical consequence is that reported causality penalties in this literature
should be read together with the evaluation protocol that produced them, and that a
real-time constraint need not be paid for in accuracy.
\end{abstract}

\begin{IEEEkeywords}
sleep staging, EEG, causal models, real-time inference, evaluation methodology
\end{IEEEkeywords}

\section{Introduction}

Most automated sleep-staging models classify 30-second epochs and are free to use signal from
later in the recording, or the whole night, when labelling an earlier epoch~\cite{deepsleepnet,
attnsleep}. That is appropriate for offline analysis of a completed study but unsuitable for a
wearable device, which must report a stage as it occurs. Causality is therefore a real design
constraint, and the natural question is what it costs.

That question is rarely answered cleanly. Published causal and non-causal models differ in
capacity, data, preprocessing and optimization simultaneously, so the gap between them is not
attributable to causality in particular. High output resolution is not the open problem either:
U-Sleep~\cite{usleep} has emitted up to 7{,}680 stages per minute since 2021, using a
bidirectional model that reads several minutes of future signal. What has not been measured is
what that resolution costs when the model may not look forward at all.

We answer it with a matched ablation: one architecture, trained twice, differing only in whether
convolutions are left-padded and attention is masked. The first result is the expected one. Under
the evaluation protocol this literature normally uses, causality costs $0.025\,\kappa$,
consistently across seeds and folds. The second result is that this number is produced by the
protocol rather than by the constraint, and this paper is about the second result.

Our contributions are threefold. We show that batch-tiled evaluation systematically penalizes
causal models, and that the penalty it reports is arithmetically the mean of a curve over
context-buffer positions, verified to within $0.003\,\kappa$. We re-measure the cost at matched
latency and find it vanishes. And we show that the apparent concentration of the penalty in REM
is the same artifact, inverting under a deployment-realistic protocol.

\section{Related Work}

Single-channel staging has moved from spectral features to end-to-end deep learning.
DeepSleepNet~\cite{deepsleepnet} pairs a dual-branch CNN with a bidirectional LSTM over epochs;
AttnSleep~\cite{attnsleep} replaces the recurrence with unmasked multi-head attention at roughly
three orders of magnitude fewer parameters. Both read future epochs. High output resolution is
solved offline: U-Time~\cite{utime} classifies every input sample and aggregates over an arbitrary
interval, and U-Sleep~\cite{usleep} evaluates that mechanism at up to 7{,}680 stages per minute
across 15{,}660 participants, reporting that finer stages carry clinical signal. It does so with a
bidirectional encoder--decoder spanning up to 13.5 minutes, roughly half of it after the labelled
instant.

A separate line targets causal or real-time operation. Bresch et al.~\cite{bresch} apply recurrent
networks to real-time single-channel staging; CausalAttenNet~\cite{causalattennet} is a causal
long-term attention network, and Micro SleepNet~\cite{microsleepnet} targets wearable hardware.
Causal dilated convolutions are known to match recurrence on long sequences while remaining
forward-looking~\cite{tcn}, and that block is used here with causally masked
attention~\cite{attention}.

To our knowledge none of this work reports a matched ablation in which one architecture is trained
with and without future access on identical folds and seeds, which is what licenses attributing a
gap to causality rather than to capacity or data. Nor, more importantly for this paper, does any of
it report the evaluation protocol used to produce a causality penalty, which we show determines the
answer.

\section{Methods}

\subsection{Data}

Sleep-EDF Expanded (sleep-cassette)~\cite{physionet, kemp}, single channel Fpz-Cz at 100~Hz, five
AASM classes~\cite{aasm}. We use the 78-subject subset (153 recordings, 1{,}629 hours after
trimming Wake to 30 minutes either side of each night's sleep period, applied per recording).
Splits are subject-wise 5-fold, verified disjoint by unit test. Expert labels are per-30-second
epoch and replicated across their 30 one-second targets, since no public dataset provides
per-second scoring; $\kappa$ is therefore the more meaningful summary and is used throughout.

\subsection{Model and the Matched Ablation}

The network maps raw EEG to one five-class prediction per second in 30{,}757 parameters: two
parallel causal convolution branches (kernels of 0.5~s and 4~s) max-pooled to 1~Hz, three residual
blocks of causal dilated convolutions (dilations 1, 2, 4), four-head masked self-attention, and a
linear classifier. Causality is enforced in exactly two places: every convolution is left-padded by
$(K-1)d$, and the attention mask sets $M_{ij} = -\infty$ for $j > i$. The non-causal arm removes
both and changes nothing else, so parameter count, width, depth, data, folds and seeds are
identical and the comparison isolates access to future signal.

A perturbation test confirms causality rather than assuming it: perturbing the input at a future
second leaves every earlier logit bit-identical, through normalization as well as the network, and
a control confirms the non-causal arm does leak. Streaming a complete held-out recording one sample
at a time through the deployment front end reproduces the offline labels exactly (79{,}440 seconds,
zero mismatches).

\begin{table}[t]
\caption{Verified Properties (Repository Unit Tests)}
\label{tab:verified}
\centering
\footnotesize
\begin{tabular}{@{}p{0.60\columnwidth}l@{}}
\toprule
Property & Result \\
\midrule
Normalization ignores future samples & PASS (bit-identical) \\
Offline arrays equal the online filter & PASS ($5.2\times10^{-14}$) \\
Raw sample $\to$ normalization $\to$ logits & PASS (bit-identical) \\
Probe does perturb later outputs & PASS (non-vacuous) \\
Non-causal control does leak & PASS \\
Incremental buffer equals offline eval. & PASS (79{,}440~s, 0 diff.) \\
Hold-out splits disjoint and complete & PASS \\
\bottomrule
\end{tabular}
\end{table}

\subsection{Normalization}

Per-epoch $z$-scoring, near-universal in this literature, normalizes an early sample using
statistics computed over its whole epoch and so reads that epoch's future. We replace it with a
trailing 30~s statistic computable one sample at a time, making the path from raw sample to logit
end-to-end causal. This is a modelling gain rather than a cost, improving both arms ($+0.024$ and
$+0.021\,\kappa$), so neither depended on the leaked information. All results below use it.

\subsection{The Two Evaluation Protocols}

\textbf{Tiled.} The recording is covered by non-overlapping context windows and every position of
every window contributes one scored prediction. This is the default in every implementation we have
examined, including our own earlier work. Each window begins from a zero-padded empty history.

\textbf{Streaming.} The model recomputes every 30~s against a full buffer and only the freshest 30
predictions are kept, so after the opening window every scored second is decided with at least
90~s of genuine history. This is what a deployed system does.

Both protocols score every second of every held-out recording exactly once, so the two are directly
comparable on identical data.

\subsection{Training and Statistics}

Focal loss ($\gamma = 2.0$) with class weights, Adam at $10^{-3}$, weight decay $10^{-4}$, batch 64,
early stopping on validation $\kappa$ with patience 8. Context is 120~s. The 5-fold protocol is
repeated under three seeds per arm, giving 15 paired seed--fold cells. Because seeds reseed the
subject partition as well as the initialization, these are fifteen distinct train/test partitions
rather than five measured three times, and repeats reuse subjects; we therefore report the
Nadeau--Bengio corrected resampled $t$-test as primary, with the naive paired statistic alongside.
Pairing is exact: within a seed both arms use byte-identical subject lists.

\section{Results}

\subsection{Headline Performance}

Table~\ref{tab:headline} states what the model achieves, so that the methodological result below
can be judged against a system of known quality. Under the streaming protocol on Sleep-EDF-78 it
reaches $\kappa = 0.700$ per second, rising to $0.707$ when majority-voted to 30~s epochs, in
30{,}757 parameters at $0.026$~ms per second of EEG on an Intel Core i9-14900HX CPU. The
subject-level bootstrap 95\% interval on tiled $\kappa$, resampling whole subjects rather than
seconds, is $[0.642, 0.686]$. This is below the strongest published single-channel models, which
read future signal and use two to three orders of magnitude more parameters; the result of this
paper concerns what separates them, and Section~\ref{sec:matched} bounds the contribution of
causality to that gap at approximately zero.

\begin{table}[t]
\caption{Model Performance (Sleep-EDF-78, 15 Folds Pooled)}
\label{tab:headline}
\centering
\begin{tabular}{lcc}
\toprule
Metric & Tiled & \textbf{Streaming} \\
\midrule
Accuracy (per second)         & 0.751 & \textbf{0.777} \\
Cohen's $\kappa$ (per second) & 0.665 & \textbf{0.700} \\
Macro-F1 (per second)         & 0.689 & \textbf{0.723} \\
$\kappa$ at 30~s              & 0.685 & \textbf{0.707} \\
\midrule
Parameters / CPU inference    & \multicolumn{2}{c}{30{,}757 / 0.026 ms per s} \\
\bottomrule
\end{tabular}
\end{table}

\subsection{The Protocol Decides the Answer}
\label{sec:matched}

Table~\ref{tab:protocol} scores the same trained weights under both protocols. They give opposite
answers.

\begin{table}[t]
\caption{Cost of Causality Under Each Evaluation Protocol\\(Sleep-EDF-78, 15 Paired Folds, Cohen's $\kappa$)}
\label{tab:protocol}
\centering
\begin{tabular}{lcccc}
\toprule
Protocol & Causal & Non-causal & $\Delta$ & Causal loses \\
\midrule
Tiled     & 0.6649 & 0.6898 & $-0.0249$ & 14/15 \\
Streaming & \textbf{0.6998} & 0.6934 & $\mathbf{+0.0065}$ & 7/15 \\
\bottomrule
\end{tabular}
\end{table}

Moving to the streaming protocol raises the causal arm by $0.0350\,\kappa$ and the non-causal arm by
$0.0036$. The per-seed gains for the causal arm are $+0.0348$, $+0.0345$ and $+0.0355$, so this is
not seed variation. Under the tiled protocol the causal penalty is significant (Nadeau--Bengio
$t = -4.76$, $p = 0.0003$; naive paired $t(14) = -10.37$, $p = 5.9\times10^{-8}$). Under the
streaming protocol the causal arm is ahead by $0.0065$ and ahead in 8 of 15 folds, which we read as
a tie rather than a win.

\subsection{Why: the Tiled Score Is a Mean Over Buffer Positions}

Define permitted latency $L$: the label for second $t$ is emitted at $t+L$, so the model sees $L$
seconds of future and $120 - L$ seconds of past. $L = 0$ is real time. A single forward pass over a
window yields a prediction at every position, and position $p$ carries exactly $119 - p$ seconds of
lookahead, so one sweep measures every latency.

Averaging that curve over all 120 positions reproduces the tiled score to within $0.003$ for both
arms: $0.6626$ against a measured $0.6649$ for the causal arm, and $0.6890$ against $0.6898$ for the
non-causal arm. The tiled protocol is, arithmetically, this average.

Table~\ref{tab:quartile} shows where the penalty lives. At positions with 91--120~s of history the
causal arm is ahead. The entire aggregate penalty is generated by positions with under roughly 60~s
of history, and most of it by positions with under 30~s --- positions a deployed system occupies
once, at start-up.

\begin{table}[t]
\caption{Where the Tiled Penalty Comes From (Cohen's $\kappa$)}
\label{tab:quartile}
\centering
\begin{tabular}{lccc}
\toprule
Buffer position & Causal & Non-causal & $\Delta$ \\
\midrule
$L$ 0--29\hphantom{0} (91--120~s history) & \textbf{0.6982} & 0.6930 & $-0.0052$ \\
$L$ 30--59 (61--90~s history)             & 0.6923 & 0.6969 & $+0.0046$ \\
$L$ 60--89 (31--60~s history)             & 0.6753 & 0.6911 & $+0.0158$ \\
$L$ 90--119 (1--30~s history)             & 0.5847 & 0.6752 & $+0.0904$ \\
\bottomrule
\end{tabular}
\end{table}

Fig.~\ref{fig:latency} gives the full curve. The causal arm declines monotonically from $L = 0$: it
cannot read the future, so larger $L$ only costs it history, and at $L = 119$ with one second of
history it collapses to $\kappa = 0.317$. The non-causal arm traces an inverted U, gaining lookahead
while losing history, and peaks at $\kappa = 0.698$ at $L = 37$~s. It overtakes the causal arm only
beyond $L \approx 27$~s, and at its optimum does not exceed what the causal arm achieves in real
time.

\begin{figure}[t]
\centering
\includegraphics[width=\columnwidth]{figures/fig_latency.png}
\caption{Cohen's $\kappa$ against permitted latency $L$, both arms, 5-point moving average. Larger
$L$ means less history and more lookahead within a fixed 120~s buffer. The causal arm is best at
$L = 0$. The causal collapse at large $L$ is the cold-start behaviour that the tiled protocol
averages into its score.}
\label{fig:latency}
\end{figure}

\subsection{The REM Concentration Is the Same Artifact}

Under the tiled protocol REM appears to pay $0.060$ F1, four times the next-largest class, and is
the only class surviving Holm correction across the five per-class tests. Table~\ref{tab:perclass}
shows this does not survive a change of protocol: under streaming the REM difference is $+0.002$,
and at $L = 0$ specifically the causal arm is ahead on REM by $0.025$.

\begin{table}[t]
\caption{Per-Class F1 by Protocol. Positive $\Delta$ Means the Causal Arm Pays}
\label{tab:perclass}
\centering
\begin{tabular}{lccc|ccc}
\toprule
 & \multicolumn{3}{c|}{Tiled} & \multicolumn{3}{c}{Streaming} \\
Class & Caus. & Non-c. & $\Delta$ & Caus. & Non-c. & $\Delta$ \\
\midrule
W   & 0.893 & 0.911 & $+0.018$ & 0.904 & 0.910 & $+0.006$ \\
N1  & 0.421 & 0.446 & $+0.025$ & 0.463 & 0.454 & $-0.009$ \\
N2  & 0.773 & 0.775 & $+0.002$ & 0.796 & 0.780 & $-0.017$ \\
N3  & 0.665 & 0.673 & $+0.008$ & 0.697 & 0.678 & $-0.020$ \\
REM & 0.695 & 0.754 & $\mathbf{+0.060}$ & 0.755 & 0.757 & $\mathbf{+0.002}$ \\
\bottomrule
\end{tabular}
\end{table}

\begin{figure}[t]
\centering
\includegraphics[width=\columnwidth]{figures/fig_perclass.png}
\caption{Per-class F1 the causal arm gives up, under batch tiling (averaged over all buffer
positions) and at deployment latency $L = 0$. The apparent REM penalty inverts. Values are from the
latency sweep, which scores each latency on a one-in-five subsample of seconds;
Table~\ref{tab:perclass} gives the full-coverage figures.}
\label{fig:perclass}
\end{figure}

The explanation is the same and is mechanistically sensible. REM is defined partly by surrounding
context, which is why a technician resolves it against neighbouring epochs rather than in isolation.
It is therefore the class most damaged by a nearly empty buffer, and cold-start positions are
exactly what the tiled protocol averages in. What survives is a more careful statement than the one
it replaces: under tiled evaluation the apparent penalty concentrates in REM because REM is the
class most sensitive to history, not because REM depends on future signal.

\subsection{Scope: the Fixed Buffer}

Latency here is measured within a buffer of fixed length, so lookahead is bought with history. An
offline system would enlarge the buffer and take both. For the causal arm we can settle this
exactly: masking makes its output at position $p$ a function of inputs $0 \ldots p$, so its
context-response curve is buffer-invariant, and measured in buffers of 120, 240 and 360~s it agrees
to within $0.0035$ at every history value. That curve peaks near 145~s of history
($\kappa = 0.705$), so a deployed causal system should use a slightly longer buffer than we did.

\begin{table}[t]
\caption{Causal Arm: $\kappa$ Against Available History}
\label{tab:context}
\centering
\footnotesize
\setlength{\tabcolsep}{3.4pt}
\begin{tabular}{lccccccc}
\toprule
History (s) & 30 & 60 & 90 & 120 & \textbf{145} & 240 & 360 \\
\midrule
$\kappa$ & 0.658 & 0.690 & 0.701 & 0.704 & \textbf{0.705} & 0.698 & 0.683 \\
\bottomrule
\end{tabular}
\end{table}

For the non-causal arm the same manoeuvre is not interpretable. Enlarging its buffer lengthens the
sequence its unmasked attention spans, and holding history fixed while increasing only lookahead
reduces $\kappa$ monotonically ($0.694 \to 0.671 \to 0.655$ at 60~s of history for buffers of 120,
240 and 360~s). Extra lookahead cannot itself degrade a model free to ignore it, so this is length
extrapolation rather than a property of lookahead. We therefore scope the result to the buffer both
arms were trained for, and note that a non-causal model \emph{trained} at longer context is the
experiment that would extend or overturn it.

\section{Discussion and Limitations}

An earlier version of this work reported that causality costs $0.029\,\kappa$ concentrated in REM.
Both statements reproduce under the protocol that produced them and both are artifacts of it. We
report them alongside the corrected measurement because the discrepancy is the finding.

The mechanism is not subtle once seen. A streaming model is defined by what it does with accumulated
history; tiling repeatedly destroys that history and scores the model on the wreckage, and it does so
asymmetrically, because a model that reads forward is largely unaffected at the start of a window
while one that cannot is crippled there.

We do not claim causality is free in general, and three limits bound the claim. The streaming protocol
denies the non-causal arm the lookahead that defines it, so this compares deployment options rather
than architectures in the abstract; the fair framing is matched latency, under which lookahead is
worth nothing the causal model cannot obtain from history, while at unbounded latency it remains worth
roughly $0.025\,\kappa$. The non-causal arm is our own matched twin, not a state-of-the-art offline
model, so the result bounds what causality costs within this capacity class. And the demonstration is
on one architecture, one channel and one dataset family; the mechanism is generic but its magnitude
elsewhere is unknown. Supervision is also 30-second labels replicated to 1~Hz, and no inter-scorer
ceiling exists for Sleep-EDF, so $\kappa$ is positioned against published models rather than against a
second expert.

The practical implication is narrow and useful. A system that must run in real time should not expect
to pay for it in accuracy, and should not be evaluated with a protocol that makes it appear to. We
would encourage any comparable ablation to report both protocols, since the two numbers here differ by
more than the effect being measured.

\section{Conclusion}

Measured with the batch-tiling protocol standard in this literature, a strict causal constraint costs
$0.025\,\kappa$ in single-channel sleep staging and appears concentrated in REM. Measured at matched
latency it costs nothing, and the REM concentration inverts. The tiled score is arithmetically the mean
over context-buffer positions, and the penalty it reports is generated by positions a streaming system
occupies once at start-up. Reported causality penalties should be read together with the protocol that
produced them.

\section*{Reproducibility}
\sloppy

Implementation, configurations, archived per-fold results, per-second predictions and the unit tests
covering the causality and split properties are at
\url{https://github.com/swapnil5053/causal-sleep-staging}. Every analysis reported here runs from the
archived predictions without retraining. Sleep-EDF Expanded is public~\cite{physionet,kemp} and is not
redistributed; this is a secondary analysis of de-identified public data and required no ethics review.

\section*{Acknowledgment}

The authors thank Dr.\ Bhaskarjyoti Das for supervision throughout this project.

\begin{thebibliography}{99}
\small

\bibitem{deepsleepnet} A. Supratak \emph{et al.}, ``DeepSleepNet: A model for automatic sleep stage
scoring based on raw single-channel EEG,'' \emph{IEEE Trans. Neural Syst. Rehabil. Eng.}, vol. 25,
no. 11, pp. 1998--2008, 2017.

\bibitem{attnsleep} E. Eldele \emph{et al.}, ``An attention-based deep learning approach for sleep
stage classification with single-channel EEG,'' \emph{IEEE Trans. Neural Syst. Rehabil. Eng.},
vol. 29, pp. 809--818, 2021.

\bibitem{usleep} M. Perslev \emph{et al.}, ``U-Sleep: Resilient high-frequency sleep staging,''
\emph{npj Digit. Med.}, vol. 4, art. 72, 2021.

\bibitem{physionet} A. L. Goldberger \emph{et al.}, ``PhysioBank, PhysioToolkit, and PhysioNet,''
\emph{Circulation}, vol. 101, no. 23, pp. e215--e220, 2000.

\bibitem{kemp} B. Kemp \emph{et al.}, ``Analysis of a sleep-dependent neuronal feedback loop,''
\emph{IEEE Trans. Biomed. Eng.}, vol. 47, no. 9, pp. 1185--1194, 2000.

\bibitem{aasm} R. B. Berry \emph{et al.}, \emph{The AASM Manual for the Scoring of Sleep and
Associated Events, Version 2.6}. Darien, IL: AASM, 2020.

\bibitem{utime} M. Perslev \emph{et al.}, ``U-Time: A fully convolutional network for time series
segmentation applied to sleep staging,'' in \emph{NeurIPS}, 2019.

\bibitem{bresch} E. Bresch \emph{et al.}, ``Recurrent deep neural networks for real-time sleep
stage classification from single channel EEG,'' \emph{Front. Comput. Neurosci.}, vol. 12, art. 85,
2018.

\bibitem{tcn} S. Bai, J. Z. Kolter, and V. Koltun, ``An empirical evaluation of generic convolutional
and recurrent networks for sequence modeling,'' \emph{arXiv:1803.01271}, 2018.

\bibitem{attention} A. Vaswani \emph{et al.}, ``Attention is all you need,'' in \emph{NeurIPS}, 2017,
pp. 5998--6008.

\bibitem{focal} T.-Y. Lin \emph{et al.}, ``Focal loss for dense object detection,'' in \emph{ICCV},
2017, pp. 2980--2988.

\bibitem{nadeau} C. Nadeau and Y. Bengio, ``Inference for the generalization error,'' \emph{Mach.
Learn.}, vol. 52, no. 3, pp. 239--281, 2003.

\bibitem{causalattennet} J. Pan \emph{et al.}, ``CausalAttenNet: A fast and long-term-temporal network
for automatic sleep staging with single-channel EEG,'' \emph{IEEE Trans. Instrum. Meas.}, vol. 73,
pp. 1--13, 2024.

\bibitem{microsleepnet} G. Liu \emph{et al.}, ``Micro SleepNet: Efficient deep learning model for
mobile terminal real-time sleep staging,'' \emph{Front. Neurosci.}, vol. 17, art. 1218072, 2023.

\end{thebibliography}

\end{document}
